% Target: rmp-iii-b-braid-four
% Formula: The final braid places R_{i,b} above R_{j,a} to resolve two crossings.
% Ink: Kernel plane; two labelled four-port resolver boxes, two parallel
%   inter-resolver strings, an upper return arch, and the lower-rail bead.
% Boundary: Two horizontal rails and the diagonal transversal are open; red
%   continuations exit at the middle left and below the lower rail.
% Source: paper TN-Review-main.tex line 1680; author source ImagesReview Section
%   III/ImagesReview Section III.tex lines 900-924
% Capabilities: kernel, plane-frame, braid-resolver, crossing-order
\[
  \tenkzkernel
  \tndeclare{species}{operator}{hue=source:red}
  \begin{tenkz}[
      rows={wire,wire,wire,wire,wire,wire}, cols=6,
      frame=plane, bonds=none]
    % Anyons i and j sit where the diagonal transversal meets the two rails.
    \tn[at=(1,2), name=i, role=marked, label pos=nw]{i}
    \tn[at=(6,2), name=j, role=extra, label pos=sw]{j}
    % the stacked resolver boxes; drawn over the strings, they conceal the
    % resolved order exactly as the author's white boxes do
    \tn[at=(2,3), skin=box,
        name=Rib, wide=2, wires=2,
        ports={90@1:virtual, 90@2:virtual, 270@1:virtual, 270@2:virtual}]{R_{i,b}}
    \tn[at=(4,3), skin=box,
        name=Rja, wide=2, wires=2,
        ports={90@1:virtual, 90@2:virtual, 270@1:virtual, 270@2:virtual}]{R_{j,a}}
    % two rails and the transversal line through both anyons
    \tnwire{open w}{i.w} \tnwire{i.e}{open e}
    \tnwire{open w}{j.w} \tnwire{j.e}{open e}
    \tnwire{open n}{i.n} \tnwire{i.s}{j.n} \tnwire{j.s}{open s}
    % Resolver attachments use declared ports. The two branches
    % between resolvers stay separate, as in the author's final braid.
    \tnwire[kind=string, species=operator, name=si]
            {i.315}{Rib.90@2}
    \tnwire[kind=string, species=operator, name=right]
            {Rib.270@2}{Rja.90@2}
    \tnwire[kind=string, species=operator, name=tail, via={xa}]
            {Rja.270@2}{open s}
    \tnwire[kind=string, species=operator, name=sj, route=arc]
            {j.20}{Rja.270@1}
    \tnwire[kind=string, species=operator, name=left]
            {Rja.90@1}{Rib.270@1}
    \tnwire[kind=string, species=operator, name=return,
            via={0.1 n of Rib.90@1, 0.2 ne of xj}]
            {Rib.90@1}{xj.45}
    \tnwire[kind=string, species=operator, name=west]
            {xj.180}{open w}
    % marker beads where strings pierce lattice lines
    \tn[at=0.2 e of (6,4), role=marked, name=xa]{}
    \tn[at=midway i and j, role=extra, name=xj]{}
  \end{tenkz}
\]
